% =====================================================================
%  SUBMISSION PACKAGE v2 — RENEWABLE ENERGY (Elsevier), elsarticle
%
%  Compile:  pdflatex RE_manuscript_v2.tex  (x2)
%  Submission form: \documentclass[review]{elsarticle} gives double
%  spacing + line numbers, which Elsevier review expects.
%
%  STATUS MARKERS
%    [NEW]   written in full in this version
%    [CARRY] prose exists in draft v2 (.docx); port it, do not retype
%            numbers — every table below is transferred verbatim from
%            the .docx extraction
%    [TO COMPUTE] requires a run not yet made
%
%  RULE: no number appears here unless it was transferred from a draft
%  v2 table or is reproduced by a script in the code package.
% =====================================================================

\documentclass[review,3p,times]{elsarticle}

\usepackage{amsmath,amssymb}
\usepackage{booktabs,array,longtable,graphicx}
\usepackage[hidelinks]{hyperref}
\usepackage{url}
\usepackage{textcomp}
\journal{Renewable Energy}

\newcommand{\todo}[1]{\textbf{[TO COMPUTE: #1]}}
\newcommand{\carry}[1]{\textsf{\small [CARRY --- #1]}}

\begin{document}

% =====================================================================
\begin{frontmatter}

\title{Leakage-free short-term wind power forecasting at turbine level:
a low-latency causal filter bank, a training-free causality certificate,
and the cost of honest evaluation}

% ---- authors: order to be fixed before submission; see note in README
\author[tog-ce]{Gökhan Gülmez}
\author[tog-ie]{Kadir Sarıkaya\corref{cor1}}
\cortext[cor1]{Corresponding author.}
\ead{kadir.sarikaya@gop.edu.tr}

\affiliation[tog-ce]{organization={Department of Computer Engineering,
  Faculty of Engineering and Architecture, Tokat Gaziosmanpaşa University},
  city={Tokat}, country={Türkiye}}
\affiliation[tog-ie]{organization={Department of Industrial Engineering,
  Faculty of Engineering and Architecture, Tokat Gaziosmanpaşa University},
  city={Tokat}, country={Türkiye}}

% ---------------------------------------------------------------
% HIGHLIGHTS — Elsevier: 3–5 bullets, each at most 85 characters
% including spaces. Counts given in comments; verify in the
% submission system, which enforces the limit.
% ---------------------------------------------------------------
\begin{highlights}
\item Full-series VMD understates causal MAE by 34--57\% on six SDWPF turbines % 70
\item A training-free certificate separates causal from leaky decomposition % 74
\item A locked causal filter bank matches rolling VMD at 1300x lower latency % 75
\item Under a causal protocol decomposition adds nothing over raw SCADA inputs % 76
\item An oracle-normalised ratio shows how much admissible leak is realised % 73
\end{highlights}

% ---------------------------------------------------------------
% ABSTRACT  [NEW] — 248 words. Reframed for Renewable Energy: the
% operational consequence (deployable accuracy, update cost) leads.
% ---------------------------------------------------------------
\begin{abstract}
Decomposition-based hybrids dominate the short-term wind power forecasting
literature, but they are usually evaluated by decomposing the whole power
series before the train/test split, so the features describing a test-period
instant encode that instant's own future. Honest evaluation requires
recomputing the decomposition from past data only at every forecast origin,
and that is computationally expensive --- which is why it is rarely done. This
study quantifies the artefact and removes the cost barrier for turbine-level
SCADA forecasting. On six turbines of the SDWPF farm, over horizons from
0.5 to 6~h and eight random initialisations, full-series variational mode
decomposition (VMD) understates the mean absolute error of the identical
causal pipeline by 34--57\%; the gap is significant on every turbine under
Diebold--Mariano tests with Holm correction and is unchanged when the
forecaster is trained with an absolute-error loss. A look-ahead dose--response
experiment attributes the gap to future information rather than to amplitude,
phase or boundary effects. We introduce two instruments: a training-free
causality certificate that probes an operator with future perturbations and
prefix truncation, and reports a leakage reach and gain rather than a pass/fail
verdict; and an oracle-normalised ratio that expresses how much of the leakage
an operator admits it actually realises. We then show that a causal spectral
filter bank whose bands are locked from the training period attains causal
accuracy statistically indistinguishable from rolling VMD at roughly
1300 times lower per-update latency, making leakage-free evaluation affordable
at farm scale. Under the causal protocol, decomposition yields no consistent
accuracy gain over raw SCADA inputs.
\end{abstract}

\begin{keyword}
wind power forecasting \sep SCADA \sep temporal data leakage \sep
variational mode decomposition \sep causal preprocessing \sep
evaluation protocol \sep Diebold--Mariano test
\end{keyword}

\end{frontmatter}

% =====================================================================
\section{Introduction}\label{sec:intro}
% =====================================================================
% [NEW]

\subsection{Operational motivation}

Short-term wind power forecasts drive unit commitment, reserve procurement,
intraday trading and curtailment decisions, so the accuracy figure a
forecaster reports is not an academic quantity: it is the basis on which an
operator sizes reserve. A figure that cannot be reproduced in deployment
therefore mis-sizes reserve, and the error is one-sided --- published accuracy
that is too good leads to reserve that is too small.

Among published short-term wind power forecasters, decomposition-based
hybrids are the most common design: the observed power series is separated
into frequency components and a learner is trained on those components.
Qian et al.~\cite{qian2019} survey the family. The design rests on an
evaluation choice that is rarely stated: \emph{when} the decomposition is
fitted. Variational mode decomposition (VMD)~\cite{dragomiretskiy2014} and the
empirical-mode family are global, non-causal transforms. If one is fitted once
over a series that contains the test period, the component value at time $t$
is a function of observations after $t$. The target is never shown to the
model, but information about it is.

\subsection{Why the honest protocol is not adopted}

The remedy is known: recompute the decomposition from past-only data at each
forecast origin. Chen et al.~\cite{chen2022} catalogue the problem as the
``boundary issue'' and classify solutions as algorithm-, sampling- or
iteration-based. The obstacle is cost. Rolling redecomposition requires a new
variational optimisation at every forecast instant; Jeon et al.~\cite{jeon2026}
state plainly that this carries ``a real computational cost'' and call it ``the
price of an evaluation that reflects deployment''. At a farm of the size of
SDWPF (134 turbines) that price is prohibitive, and the literature has
responded by not paying it.

This paper takes the cost seriously as an engineering problem rather than as a
caveat. Its central practical claim is that the price can be reduced by about
three orders of magnitude without measurable loss of causal accuracy, which
makes the honest protocol usable at farm scale.

\subsection{Position of this work}\label{sec:position}

We do not claim to be the first to observe or to measure the artefact, and the
reader should be clear about what is already established. Rolling or
sample-wise decomposition has been proposed as a remedy for gasoline
consumption~\cite{yu2021}, crude oil prices~\cite{xu2023}, power
load~\cite{xuA2025,wangj2024}, wind speed~\cite{wang2024,yangD2024,tan2025,pang2026,zhou2024}
and ultra-short-term wind power~\cite{huo2024}. Quantified comparisons also
exist: Xu and Wang~\cite{xu2023} contrast rolling and whole-series EEMD under
the same predictors; Zhang et al.~\cite{zhang2025} report, across six datasets
and three decompositions, that the abnormal error reduction is carried mainly
by the high-frequency components; Zhou et al.~\cite{zhou2024} quantify the
overestimation in wind speed forecasting as 40--55\% in mean absolute and
65--85\% in mean squared error; and Jeon et al.~\cite{jeon2026} report, for
German load, solar and wind generation plus a benchmark series, that the
whole-series protocol understates the error of the identical causal pipeline
by 30--93\%, attributing the artefact to look-ahead rather than to history
length. Valdivieso et al.~\cite{valdivieso2026} report a comparable magnitude
for load.

Three things are left open by that body of work, and they are what this paper
contributes.

\emph{First, the cost of the remedy is named but not addressed.} None of the
studies above measures, or reduces, the per-update latency of causal
decomposition. Section~\ref{sec:lcsf} defines a locked causal spectral filter
bank and Section~\ref{sec:cost} measures its latency against rolling VMD on
the same series.

\emph{Second, there is no operational test of an operator's causality.}
Jeon et al.~\cite{jeon2026} offer a five-point checklist, and Kapoor and
Narayanan~\cite{kapoor2023} offer reporting sheets; neither can be applied to
a given operator to decide whether it leaks. Section~\ref{sec:certificate}
defines a training-free certificate that can, and reports a profile rather
than a verdict.

\emph{Third, the magnitude has been reported only as a raw error gap.} A raw
gap is divided by the causal error, so it grows when the causal pipeline is
weak and operators with different ceilings are not comparable.
Section~\ref{sec:lambda} anchors the gap to an attainable, operator-relative
oracle.

The setting is also different in a way that matters for this journal:
all the measurements above are on aggregate or national series at hourly
resolution, whereas reserve and curtailment decisions are taken on
turbine-level SCADA at 10--30~min. Turbine-level power has no spatial
smoothing, much higher volatility, and pitch and curtailment artefacts. We
therefore report the measurement at turbine level, under a single causal
protocol, across turbines, horizons, random initialisations and forecaster
families.

\subsection{Contributions}

\begin{enumerate}
\item A \textbf{locked causal spectral filter bank (LCSF)} whose bands are
learned from the training period only, attaining causal accuracy
statistically indistinguishable from rolling VMD at roughly 1300 times lower
per-update filtering latency (Sections~\ref{sec:lcsf}, \ref{sec:cost}).
It is not offered as a more accurate forecaster; its value is that it makes
the honest protocol affordable.
\item A \textbf{training-free causality certificate}: a future-perturbation
sensitivity profile $D(\tau;\delta)$ plus a prefix-invariance test, reporting
a leakage \emph{reach} and \emph{gain}. It distinguishes leaky operators from
one another --- the sensitivity of full-series VMD does not decay with lag,
whereas that of two-sided filtering does (Sections~\ref{sec:certificate},
\ref{sec:certres}). Resolved per component and weighted by the forecaster's
reliance on each component, it yields a \textbf{weighted leakage exposure} and
a graded acceptance criterion calibrated against run-to-run variability, in
place of a pass/fail verdict on exact zero
(Sections~\ref{sec:exposure}, \ref{sec:exposureres}).
\item An \textbf{oracle-normalised leakage ratio} $\lambda$, bounded on an
attainable operator-relative oracle, which orders the two leaky operators
opposite to the raw gap (Sections~\ref{sec:lambda}, \ref{sec:lambdares}).
\item A \textbf{turbine-level measurement} of the leakage-associated gap with
Diebold--Mariano testing across six turbines, six horizons, eight seeds and
four forecaster families, including the native 10-min grid.
\item \textbf{Negative results reported in full}: under the causal protocol
decomposition gives no consistent gain over raw SCADA inputs, and the position
of causal models against persistence depends on the training loss.
\end{enumerate}

% =====================================================================
\section{Related work}\label{sec:related}
% =====================================================================
% [NEW structure; absorbs draft v2 §2.1–§2.2, reordered so that the
% closest prior art leads]

\subsection{Decomposition-based wind power forecasting}
\carry{draft v2 §2.1. Keep the VMD description. Add the filter-bank reading:
Dragomiretskiy and Zosso~\cite{dragomiretskiy2014} describe VMD as a
generalisation of the Wiener filter into multiple adaptive bands; this is the
theoretical warrant for Section~\ref{sec:lcsf}. Add 8--12 recent
decomposition-hybrid wind papers --- the literature being audited must be
represented or the premise is unsupported.}

\subsection{Controlled measurements of the artefact}
\textsf{[NEW --- this subsection leads, because it contains the closest prior
art.]} Jeon et al.~\cite{jeon2026} compare five feature regimes (no
decomposition, full-series, split-then-decompose, rolling causal windows of
30--365 days, and an expanding causal window) on four hourly series with three
model families, ten seeds and Diebold--Mariano, Wilcoxon and Friedman tests.
They report a 30--93\% understatement in 11 of 12 dataset--model combinations,
attribute 91\% of the gap to look-ahead via the expanding-window regime, find
that split-then-decompose leaves 77--101\% of the artefact intact, and find
causally decomposed pipelines no better than undecomposed ones in 10 of 12
combinations. Zhang et al.~\cite{zhang2025} reach the same qualitative
conclusion over six datasets and three decompositions, and additionally
attribute the artefact to the high-frequency components.
Zhou et al.~\cite{zhou2024} quantify it for wind speed (40--55\% in PMAE,
65--85\% in PMSE) and propose a rolling-decomposition model with a time-series
extension. Valdivieso et al.~\cite{valdivieso2026} report 46--51\% for load.

\textsf{Required text:} state explicitly what this study adds to these, using
the three openings of Section~\ref{sec:position}, and do not restate the
measurement itself as novel.

\subsection{Leakage-preventing designs}
\carry{draft v2 §2.2, second half. Chen et al.~\cite{chen2022};
Yang et al.~\cite{yang2024}; Feng et al.~\cite{feng2026};
Sesay~\cite{sesay2026}; Wang et al.~\cite{wang2024};
Tan et al.~\cite{tan2025}; Haq et al.~\cite{haq2026};
Pang and Dong~\cite{pang2026}; Huo et al.~\cite{huo2024};
Wang et al.~\cite{wangj2024}; Quilty and Adamowski~\cite{quilty2018};
Maslova et al.~\cite{maslova2015}.}

\subsection{Evaluation practice in wind power forecasting}
\textsf{[NEW]} The reference-model and metric conventions used here follow
Nielsen et al.~\cite{nielsen1998}, Madsen et al.~\cite{madsen2005} and
Messner et al.~\cite{messner2020}. González Sopeña et al.\ review performance
evaluation methods for short-term statistical wind power
forecasting~\cite{sopena2020} and propose a benchmarking framework applied to
mode-decomposition models~\cite{sopena2022}; Lee et al.~\cite{lee2022} publish
an open-source systematic validation framework and a benchmark exercise.
Hong et al.~\cite{hong2020}, reviewing the field, advise authors who obtain
exaggerated results to check ``if future information has leaked into the
process''. Deignan et al.~\cite{deignan2025} study the portability and
hyperparameter sensitivity of VMD- and EEMD-based wind power forecasters on
SCADA data from Ireland and the UK and report that about four modes suffice.
General treatments of leakage~\cite{kaufman2012,kapoor2023,rosenblatt2024} and
of forecast-evaluation pitfalls~\cite{hewamalage2023} complete the picture.

\subsection{Causal and online decomposition}
\textsf{[NEW]} Short-time VMD~\cite{jia2026} and block-wise recursive sliding
VMD~\cite{dan2023} reduce the cost of windowed decomposition; generalised VMD
with prior centre frequencies~\cite{guo2021} fixes the bands, as LCSF does,
but is not framed as a leakage remedy and is not benchmarked for latency in a
forecasting pipeline. \todo{benchmark LCSF against a strictly causal
short-time VMD and against recursive sliding VMD; verify the causality of the
latter from its full text before use}

\subsection{Positioning}
\carry{draft v2 Table 1, rewritten. Rows: Jeon et al.~\cite{jeon2026},
Zhang et al.~\cite{zhang2025}, Zhou et al.~\cite{zhou2024},
Huo et al.~\cite{huo2024}, this study. Columns: series and resolution;
what is compared; normalisation of the reported gap; sites / horizons / seeds;
significance test; \emph{operational instrument offered}. The last column is
where this study separates: a measured low-latency causal operator and an
applicable certificate, against checklists elsewhere.}

% =====================================================================
\section{Data and protocol}\label{sec:data}
% =====================================================================
\carry{draft v2 §3.1--3.4 in full.}
\begin{itemize}
\item \S3.1 SDWPF~\cite{zhou2024sdwpf}; Sinovel SL1500/82, 1500~kW rating; the
sampling-interval finding (15~min in 2020 and 10~min in 2021 in the file used)
and the 30-min common grid. \textbf{Required before submission:} verify this
against the figshare original and state the outcome.
\item \S3.2 windowing, horizons, persistence. \textbf{Add} the reference
models of Section~\ref{sec:reference} here.
\item \S3.3 the 37 input features.
\item \S3.4 chronological 70/15/15 split, 12-step purge, train-only scaling.
\end{itemize}
\todo{rolling-origin evaluation, or at minimum the main results by monthly
test sub-period}
\todo{external validation on a second farm (Kelmarsh, 6 turbines, 10-min
SCADA~\cite{kelmarsh2022}), chosen to match the SDWPF subset size and to avoid
the grid-harmonisation issue}

% =====================================================================
\section{Decomposition operators}\label{sec:operators}
% =====================================================================
\carry{draft v2 §4.1--4.2 and §4.4: full-series VMD (leaky), rolling VMD
(causal), and the forecaster definitions (persistence, raw LSTM, CatBoost
Direct, the LSTM variants, the residual CatBoost stage).}

\subsection{Locked causal spectral filter bank}\label{sec:lcsf}
\carry{draft v2 §4.3.} VMD is applied to the training partition only to obtain
$K$ centre frequencies; the geometric means of adjacent centre frequencies
become band edges, and a bank of third-order Butterworth filters (one
low-pass, two band-pass, one high-pass) is applied in a strictly
forward direction, with the initial state derived from the first observation.
Applying the same bands zero-phase gives the leaky counterpart, which differs
from the causal one only in filtering direction and therefore serves as a
control operator.

\textbf{Required additions.} (i) Ground the construction in the filter-bank
reading of VMD~\cite{dragomiretskiy2014} and differentiate it from generalised
VMD with prior centre frequencies~\cite{guo2021}. (ii) State that LCSF is not
presented as a mathematical equivalent of VMD at test time. (iii) Report the
second-order-section check (draft v2 §6.17): the largest difference between the
transfer-function and second-order-section implementations was at most
$8\times10^{-6}$~kW per component, so the numerical form does not affect the
results.

% =====================================================================
\section{Evaluation framework}\label{sec:eval}
% =====================================================================
% [NEW] — the methodological core of the paper

\subsection{Error metrics and significance}
\carry{draft v2 §5: MAE, RMSE, nMAE, nRMSE, $R^2$; skill relative to a
reference; Diebold--Mariano~\cite{dm1995} with Bartlett long-run variance and
the Harvey--Leybourne--Newbold small-sample correction~\cite{hln1997}; Holm
correction~\cite{holm1979}; block bootstrap~\cite{kunsch1989}.}

\subsection{Reference models}\label{sec:reference}
\textsf{[NEW]} Persistence alone is not the community reference. We adopt the
MSE-optimal convex combination of Nielsen et al.~\cite{nielsen1998},
\begin{equation}
\hat p(t+k) = a_k\,p_t + (1-a_k)\,\bar p,
\qquad
a_k = \frac{\frac{1}{N}\sum_t \tilde p_t \tilde p_{t+k}}
            {\frac{1}{N}\sum_t \tilde p_t^{\,2}},
\qquad \tilde p_t = p_t-\bar p,
\label{eq:nielsen}
\end{equation}
with the lag-$k$ autocorrelation $a_k$ and the mean $\bar p$ estimated on the
training period only and separately per turbine, as those authors advise.
As $a_k\to1$ it reduces to persistence and as $a_k\to0$ to climatology.
Because the combination minimises squared error, its MAE can exceed that of
persistence; both metrics are therefore reported, and skill is defined
following Madsen et al.~\cite{madsen2005} with the metric set of
Messner et al.~\cite{messner2020}.

\subsection{Leakage ratios}\label{sec:lambda}
For matched causal and leaky versions of one operator, the raw gap is
\begin{equation}
\mathrm{LSI} =
\frac{E_{\mathrm{causal}}-E_{\mathrm{leaky}}}{E_{\mathrm{causal}}}\times 100 ,
\label{eq:lsi}
\end{equation}
which divides by the causal error and so is inflated when the causal pipeline
is weak --- the regime a leakage study is likely to be in --- and is not
comparable across operators with different ceilings.
\textsf{[NEW]} We therefore also report an oracle-normalised ratio. Let
$E_{\mathrm{oracle}}$ be the error of the \emph{same} operator applied over a
window that contains the target, i.e.\ the best this operator could do if
allowed to see the answer. Then
\begin{equation}
\lambda =
\frac{E_{\mathrm{causal}}-E_{\mathrm{leaky}}}
     {E_{\mathrm{causal}}-E_{\mathrm{oracle}}}
\label{eq:lambda}
\end{equation}
answers: of the leakage this operator admits, what fraction does the audited
protocol realise? Oracle anchoring is established in forecast
verification~\cite{wheatcroft2019}, with the caveat that a perfect anchor is
often unattainable; Eq.~\eqref{eq:lambda} uses an attainable,
operator-relative oracle instead. Two limits must be stated with every
reported value: $\lambda$ is undefined when
$E_{\mathrm{causal}}\le E_{\mathrm{oracle}}$, and because the look-ahead curve
of rolling VMD is not monotone (Section~\ref{sec:mechanism}),
$E_{\mathrm{oracle}}$ is a defined reference point rather than a strict upper
bound. $\lambda$ complements the raw gap; it does not replace it.

\subsection{Causality certificate}\label{sec:certificate}
\textsf{[NEW]} Let $\Phi$ map a series $x$ to $K$ components, with
$\Phi(x)_{t,k}$ the $k$th component at time $t$, and let
$\mathcal F_t=\sigma(x_s: s\le t)$. The operator is \emph{causal} if
$\Phi(x)_{t,k}$ is $\mathcal F_t$-measurable for every $t$ and $k$.

Causality is a property of the computation, and for pipelines whose data
availability does not depend on data \emph{values} it can be decided by static
analysis. Full-series VMD, rolling VMD and the EMD family are not of that
kind: each is an iterative optimisation whose output at $t$ depends on the
values of the whole input window. For that class an empirical, operator-
agnostic probe is the available instrument, and we state what it establishes:
a necessary condition over a probed set, never the absence of leakage.

\paragraph{Test 1: future-perturbation profile.}
For a perturbation epoch $t_0$, magnitude $\delta$ and lag $\tau\ge0$,
\begin{equation}
D_\Phi(\tau;\delta)=
\max_{t\le t_0-\tau}\;\max_k\;
\bigl|\Phi(x)_{t,k}-\Phi\!\left(x\oplus\delta\mathbf 1_{\{s>t_0\}}\right)_{t,k}\bigr| ,
\label{eq:profile}
\end{equation}
where $x\oplus\delta\mathbf 1_{\{s>t_0\}}$ adds an independent perturbation of
scale $\delta$ to every observation strictly after $t_0$. If $\Phi$ is causal,
$D_\Phi\equiv0$. Reporting the profile rather than a single value yields the
\emph{leakage reach}, the largest probed $\tau$ at which $D_\Phi>0$, and the
\emph{leakage gain}, the growth of $D_\Phi$ in $\delta$.

\paragraph{Test 2: prefix invariance.}
Independently, require $\Phi(x_{1:t})_{t,k}=\Phi(x_{1:N})_{t,k}$ at a set of
random cut points. Test~1 probes sensitivity to a perturbation of the future;
Test~2 probes dependence on the existence of the future.

\paragraph{What it is not.} Passing does not prove causality; it establishes
that no dependence was detected at the probed epochs, lags and magnitudes, and
must be reported with those parameters. Related instruments operate on code or
on process rather than on the temporal behaviour of a data-adaptive
operator~\cite{kapoor2023,kaufman2012}.

\subsection{Weighted leakage exposure and a graded criterion}\label{sec:exposure}
% [NEW — TO COMPUTE]

Equation~\eqref{eq:profile} takes a maximum over components, so it answers only
whether an operator leaks. Two further questions matter in practice: \emph{which}
components carry the leak, and whether the forecaster actually uses them. The
first has been studied: Zhang et al.~\cite{zhang2025} add intrinsic mode
functions to a forecaster one at a time and report that the high-frequency
components carry most of the artefact. That measure, however, fuses two distinct
quantities --- how much future the operator exposes in a component, and how much
the model exploits it --- because it is read off the change in forecast error.
The certificate separates them, because $D_k$ is a property of $\Phi$ alone.

Resolve Eq.~\eqref{eq:profile} per component,
\begin{equation}
D_{\Phi,k}(\tau;\delta)=
\max_{t\le t_0-\tau}
\bigl|\Phi(x)_{t,k}-\Phi\!\left(x\oplus\delta\mathbf 1_{\{s>t_0\}}\right)_{t,k}\bigr| ,
\label{eq:profilek}
\end{equation}
and let $w_k$ be the forecaster's reliance on component $k$, normalised to
$\hat w_k = w_k/\sum_j w_j$ (we report all three measures already computed in
Appendix~C: prediction-value change, permutation $\Delta$MAE and mean absolute
SHAP). Scaling each component by its own standard deviation $\sigma_k$ to make
the sum dimensionless, the \emph{weighted leakage exposure} is
\begin{equation}
\Lambda(\tau;\delta)=\sum_{k=1}^{K}\hat w_k\,\frac{D_{\Phi,k}(\tau;\delta)}{\sigma_k}.
\label{eq:exposure}
\end{equation}
$\Lambda=0$ if and only if the operator is causal or the model ignores every
leaking component. The concentration of
$\{\hat w_k D_{\Phi,k}/\sigma_k\}_{k=1}^{K}$ indicates whether the leak is
carried by one component or spread across the bank.

\paragraph{A graded acceptance criterion.} The certificate as stated above is
binary: an operator passes only if $D$ is exactly zero. That is brittle, since a
component that leaks negligibly and is never used fails on the same footing as
one that carries the whole artefact. We therefore propose accepting a pipeline
when
\begin{equation}
\Lambda(0;\delta_{\mathrm{ref}})\le c ,
\label{eq:criterion}
\end{equation}
and we calibrate $c$ on the quantity a practitioner can already measure: $c$ is
set at the exposure whose implied error inflation falls below the run-to-run
standard deviation of the same forecaster across random initialisations
(2.4--6.6~kW in these data, Table~\ref{tab:seeds} and draft v2 §6.3). The
criterion then reads: \emph{a pipeline is acceptable when the leakage it still
admits moves the error less than retraining it does.} Both $\Lambda$ and the
raw profile are reported, so a reader can apply either.

\paragraph{Scope of the claim.} $w_k$ is model-dependent, is not an integer and
is not fixed across retrainings, so Eq.~\eqref{eq:criterion} is a weighted
\emph{acceptance criterion}, not a weighted $k$-out-of-$n$ reliability model, and
it is presented as such. The symbol $\Lambda$ is distinct from the
oracle-normalised ratio $\lambda$ of Section~\ref{sec:lambda}.

\todo{(i) resolve \texttt{leakage/certificate.py::pertürbasyon\_profili} per
component --- it currently reduces over $k$ --- and re-run script 29; the cost is
dominated by full-series VMD, about 1--2 h with the default probe grid.
(ii) take $\hat w_k$ from the existing Appendix~C importance table; no retraining
is required. (iii) calibrate $c$ against the eight-seed MAE standard deviations
already reported.}

% =====================================================================
\section{Results}\label{sec:results}
% =====================================================================
% Main text keeps ten display items; everything else moves to appendices.
% All tables below transferred verbatim from draft v2.

\subsection{Magnitude of the leakage-associated gap}\label{sec:magnitude}
\carry{draft v2 §6.1--6.2 prose.}

\begin{table}[htbp]\centering
\caption{Turbine 57, multi-horizon MAE (kW) and LSI (seed 42).}
\label{tab:horizon}
\footnotesize
\begin{tabular}{lrrrrrr}
\toprule
Horizon (h) & Persistence & Raw LSTM & CatBoost Direct &
Rolling VMD (causal) & Full-series VMD (leaky) & LSI-VMD (\%) \\
\midrule
0.5 &  77.80 &  90.04 &  78.78 &  86.76 & 53.29 & 38.6 \\
1   & 108.42 & 114.55 & 106.12 & 116.11 & 64.71 & 44.3 \\
2   & 153.78 & 164.28 & 143.35 & 158.98 & 67.73 & 57.4 \\
3   & 187.89 & 211.91 & 168.05 & 191.00 & 73.43 & 61.6 \\
4   & 214.98 & 212.26 & 187.36 & 215.86 & 96.81 & 55.2 \\
6   & 255.53 & 240.69 & 215.69 & 247.92 & 88.55 & 64.3 \\
\bottomrule
\end{tabular}
\end{table}

\begin{table}[htbp]\centering
\caption{One-hour horizon, six turbines, seed 42. MAE in kW; skill relative to
persistence (\%) in parentheses. DM: Diebold--Mariano statistic, negative when
the leaky model has the lower error; all $p<10^{-15}$.}
\label{tab:sixturbines}
\scriptsize
\begin{tabular}{lrrrrrrrrrr}
\toprule
& \multicolumn{5}{c}{Causal} & \multicolumn{2}{c}{Leaky} &
\multicolumn{2}{c}{LSI (\%)} & \\
\cmidrule(lr){2-6}\cmidrule(lr){7-8}\cmidrule(lr){9-10}
Turbine & Persist. & Raw LSTM & CatBoost & Rolling VMD & LCSF &
Full VMD & LCSF 2-sided & VMD & LCSF & DM (VMD) \\
\midrule
57  & 108.42 & 114.55 (-5.7)  & 106.12 (2.1)  & 116.11 (-7.1)  & 116.29 (-7.3)  & 64.71 (40.3) & 57.34 (47.1) & 44.3 & 50.7 & $-26.56$ \\
10  & 110.58 & 131.73 (-19.1) & 116.72 (-5.6) & 128.17 (-15.9) & 127.71 (-15.5) & 69.58 (37.1) & 55.57 (49.7) & 45.7 & 56.5 & $-27.82$ \\
83  &  95.85 & 114.43 (-19.4) &  98.84 (-3.1) & 111.64 (-16.5) & 111.34 (-16.2) & 61.67 (35.7) & 57.89 (39.6) & 44.8 & 48.0 & $-24.77$ \\
120 &  95.08 &  97.32 (-2.4)  &  93.37 (1.8)  &  95.59 (-0.5)  &  96.83 (-1.8)  & 63.42 (33.3) & 57.29 (39.7) & 33.7 & 40.8 & $-18.49$ \\
61  & 107.31 & 117.78 (-9.8)  & 107.11 (0.2)  & 119.48 (-11.3) & 119.26 (-11.1) & 70.15 (34.6) & 56.01 (47.8) & 41.3 & 53.0 & $-26.12$ \\
84  & 103.19 & 109.12 (-5.7)  & 102.40 (0.8)  & 116.34 (-12.7) & 110.15 (-6.7)  & 72.61 (29.6) & 55.03 (46.7) & 37.6 & 50.0 & $-24.66$ \\
\bottomrule
\end{tabular}
\end{table}

\begin{table}[htbp]\centering
\caption{LSI (\%), mean $\pm$ standard deviation over eight seeds (42--49).
The turbines are on one farm and are not independent replicates.}
\label{tab:seeds}
\footnotesize
\begin{tabular}{lrr}
\toprule
Turbine & LSI-VMD & LSI-LCSF \\
\midrule
57  & $39.9\pm4.4$ & $48.7\pm5.0$ \\
10  & $45.1\pm3.1$ & $53.1\pm3.8$ \\
83  & $45.2\pm2.9$ & $52.0\pm4.5$ \\
120 & $37.7\pm5.2$ & $47.0\pm5.2$ \\
61  & $42.8\pm5.7$ & $50.4\pm4.9$ \\
84  & $40.7\pm4.7$ & $50.3\pm2.9$ \\
\midrule
Across-turbine mean & $41.9\pm3.0$ & $50.3\pm2.2$ \\
\bottomrule
\end{tabular}
\end{table}

\subsection{Mechanism: the look-ahead dose--response}\label{sec:mechanism}
\carry{draft v2 §6.10--6.11, condensed. This is the strongest mechanistic
evidence and should lead the subsection; the non-monotonicity at $m=8$ and the
Ridge control are compressed to one paragraph plus Appendix~A.}

\begin{table}[htbp]\centering
\caption{Turbine 57, 1~h, three seeds: MAE (seed mean) and LSI-LCSF
(mean $\pm$ std) by the amount of look-ahead $m$ granted to the component.
Conditions with $m\ge H=2$ contain the target and are target-inclusive stress
tests, not realistic use.}
\label{tab:dose}
\footnotesize
\begin{tabular}{lrrr}
\toprule
Look-ahead $m$ (steps) & Future (min) & MAE (kW) & LSI-LCSF (\%) \\
\midrule
Causal LCSF (one-sided, reference) & 0 & 115.94 & --- \\
0 (zero-phase, no future)          & 0 & 113.27 & $2.3\pm1.4$ \\
1                                  & 30 & 85.86 & $26.0\pm1.8$ \\
2 ($=H$, target in component)      & 60 & 11.05 & $90.5\pm1.0$ \\
4                                  & 120 & 29.10 & $74.9\pm0.6$ \\
8                                  & 240 & 59.43 & $48.7\pm6.3$ \\
16                                 & 480 & 63.50 & $45.2\pm2.7$ \\
48                                 & 1440 & 61.60 & $46.8\pm5.1$ \\
$\infty$ (full series, two-sided)  & --- & 61.73 & $46.7\pm7.7$ \\
\bottomrule
\end{tabular}
\end{table}

Zero-phase filtering without look-ahead raises the gap only to 2.3\%, so the
difference between two-sided and one-sided filtering is not an artefact of
amplitude or phase. One step of look-ahead raises it to 26.0\%, and including
the target raises it to 90.5\%. A backward-shifted causal control, which moves
the component away from the window end without granting any future, changes MAE
by about 1\% (draft v2 §6.11), so the boundary effect is small in these data.

\subsection{Causality certificate}\label{sec:certres}
\carry{draft v2 §6.18 prose.} The probe set was
$\tau\in\{0,1,2,4,8,16,32\}$ steps, $\delta\in\{50,100,200,400\}$~kW,
20 perturbation epochs (5 for full-series VMD) and 50 random cut points
(10 for full-series VMD).

\begin{table}[htbp]\centering
\caption{Causality certificate, Turbine 57. $D(\tau;\delta=200\,\mathrm{kW})$
in kW. Reach is the largest probed $\tau$ with $D>0$; 32 steps is the grid
bound.}
\label{tab:cert}
\footnotesize
\begin{tabular}{lrrrlrl}
\toprule
Operator & $D(\tau{=}0)$ & $D(\tau{=}8)$ & $D(\tau{=}32)$ & Reach &
Prefix max.\ diff.\ (kW) & Prefix violations \\
\midrule
LCSF, causal              &   0.00 &   0.00 &  0.00 & none      & 0 & 0 / 50 \\
Rolling VMD, causal       &   0.00 &   0.00 &  0.00 & none      & 0 & 0 / 50 \\
LCSF, two-sided (leaky)   &  87.21 &  55.39 & 10.17 & $\ge32$   & 1{,}060 & 50 / 50 \\
Full-series VMD (leaky)   & 113.49 & 235.44 & 96.57 & $\ge32$   & 386 & 10 / 10 \\
\bottomrule
\end{tabular}
\end{table}

Both causal operators give exactly zero at every probed $\tau$ and $\delta$ and
show no prefix violation. The two leaky operators differ qualitatively: the
sensitivity of two-sided filtering decays with lag, as a finite-support
operator must, whereas that of full-series VMD does not, which is consistent
with a global variational optimisation that redistributes the whole
decomposition when any part of the series changes. The profile therefore
distinguishes \emph{kinds} of leakage, not only its presence.

\subsection{Weighted leakage exposure}\label{sec:exposureres}
% [TO COMPUTE]
\todo{Table: per operator, the per-component $D_{\Phi,k}(0;200\,\mathrm{kW})$
from Eq.~\eqref{eq:profilek}, the normalised reliance $\hat w_k$ from
Appendix~C, the resulting exposure $\Lambda$ from Eq.~\eqref{eq:exposure}, and
the concentration of the per-component products. Then the calibrated threshold
$c$ of Eq.~\eqref{eq:criterion} and which operators fall on each side of it.
Expected reading, to be confirmed rather than assumed: the two causal operators
give $\Lambda=0$ exactly; two-sided LCSF and full-series VMD are separated not
only in magnitude but in concentration, since the VMD profile does not decay
with lag whereas the filter-bank profile does
(Table~\ref{tab:cert}).}

\subsection{Oracle-normalised ratio}\label{sec:lambdares}
\carry{draft v2 §6.14 prose, including its caveats.}

\begin{table}[htbp]\centering
\caption{Turbine 57, 1~h, three seeds: the four anchors of the error scale and
$\lambda$. $E_{\mathrm{oracle}}$ is the $m=H=2$ target-inclusive condition.
All errors in kW.}
\label{tab:lambda}
\footnotesize
\begin{tabular}{lrrrrr}
\toprule
Operator & $E_{\mathrm{causal}}$ & $E_{\mathrm{leaky}}$ &
$E_{\mathrm{oracle}}$ & LSI (\%) & $\lambda$ \\
\midrule
Two-sided LCSF  & 115.94 & 61.73 & 11.05 & 46.7 & 0.52 \\
Full-series VMD & 116.78 & 68.66 & 64.64 & 41.1 & 0.92 \\
\bottomrule
\end{tabular}
\end{table}

The two measures order the operators oppositely: the raw gap is larger for
two-sided LCSF, whereas $\lambda$ shows that full-series VMD realises almost
all of the leakage its operator admits while two-sided filtering realises about
half. The raw gap compares the operators against different ceilings.
This is one turbine, one horizon and three seeds; the oracle of full-series VMD
is weak (64.64 against 68.66~kW), so $\lambda$ is sensitive in that cell.
\todo{$\lambda$ for the remaining turbines and horizons}

\subsection{Cost of the causal protocol}\label{sec:cost}
\carry{draft v2 §6.4 prose.}

\begin{table}[htbp]\centering
\caption{Forecast-instant latency, Turbine 57, single core; median of 200
positions along the series for rolling VMD.}
\label{tab:cost}
\footnotesize
\begin{tabular}{lrr}
\toprule
Method & Single forecast instant & Test set (5{,}252 points) \\
\midrule
Rolling VMD ($W=672$) & 89.9 ms & $\approx$ 8 min \\
LCSF (online update)  & 69 \textmu s & 0.36 s \\
\bottomrule
\end{tabular}
\end{table}

Causal accuracy is statistically indistinguishable at 1 and 6~h, with a small
but significant LCSF advantage at 3~h. The ratio is implementation- and
hardware-dependent; the order of magnitude ($10^3$) is the result. At farm
scale the consequence is direct: rolling VMD requires on the order of an hour
per turbine for the full series, which for 134 turbines makes per-origin
redecomposition impractical, whereas the locked bank reduces it to seconds.

\subsection{Causal models against the reference models}
\carry{draft v2 §6.15 prose, with its caveat that the Nielsen reference is the
weaker MAE reference and that positive skill against it must not be read as a
stronger result than skill against persistence.}

\begin{table}[htbp]\centering
\caption{One-hour horizon, six turbines: reference errors (kW) and MAE skill
(\%) of the causal models relative to the Nielsen reference,
Eq.~\eqref{eq:nielsen}.}
\label{tab:reference}
\scriptsize
\begin{tabular}{lrrrrrrl}
\toprule
Turbine & $a_k$ & MAE pers. & MAE Nielsen & MAE clim. & RMSE pers. &
RMSE Nielsen & Skill vs Nielsen: LSTM / VMD / LCSF / CatBoost \\
\midrule
T57  & 0.837 & 108.41 & 125.67 & 338.07 & 190.21 & 186.37 & 8.8 / 7.6 / 7.5 / 15.6 \\
T10  & 0.847 & 110.58 & 135.29 & 405.95 & 199.52 & 198.35 & 2.6 / 5.3 / 5.6 / 13.7 \\
T83  & 0.837 &  95.85 & 122.54 & 371.01 & 181.64 & 182.63 & 6.6 / 8.9 / 9.1 / 19.3 \\
T120 & 0.819 &  95.08 & 112.11 & 292.29 & 169.21 & 165.78 & 13.2 / 14.7 / 13.6 / 16.7 \\
T61  & 0.826 & 107.31 & 127.02 & 335.40 & 190.36 & 186.82 & 7.3 / 5.9 / 6.1 / 15.7 \\
T84  & 0.833 & 103.19 & 119.89 & 322.52 & 180.62 & 177.20 & 9.0 / 3.0 / 8.1 / 14.6 \\
\bottomrule
\end{tabular}
\end{table}

\subsection{Sensitivity to the training loss}
\carry{draft v2 §6.16 prose.}

\begin{table}[htbp]\centering
\caption{Turbine 57, 1~h, LSTM trained with an absolute-error loss, three
seeds (42--44). Persistence MAE is 108.41~kW.}
\label{tab:loss}
\footnotesize
\begin{tabular}{llrr}
\toprule
Model & Protocol & MAE (kW) & MAE skill vs persistence (\%) \\
\midrule
Raw LSTM              & causal              & $106.65\pm0.33$ & 1.6 \\
Rolling VMD--LSTM     & causal              & $107.22\pm0.31$ & 1.1 \\
LCSF--LSTM            & causal (one-sided)  & $107.19\pm0.76$ & 1.1 \\
Full-series VMD--LSTM & leaky               & $61.61\pm3.10$  & --- \\
LCSF--LSTM            & leaky (two-sided)   & $56.86\pm3.49$  & --- \\
\bottomrule
\end{tabular}
\end{table}

Under the matched loss the causal models reach 106.65--107.22~kW and exceed
persistence by 1.1--1.6\%, and the leakage-associated gap is unchanged
(LSI-VMD $42.5\pm2.7$\%, LSI-LCSF $46.9\pm3.4$\%). Two conclusions follow: the
main finding is not an artefact of the squared-error loss, whereas the position
of causal models relative to persistence in MAE is. The absence of a
decomposition gain also holds under the absolute-error loss --- the three causal
models stay within about 0.6~kW of one another.

\subsection{Robustness}
\carry{draft v2 §6.3, §6.7--6.9, §6.12--6.13 and §7.1's equal-training-window
control, condensed into one subsection with the detail in Appendix~B: eight
seeds; CatBoost, GRU and causal TCN; rolling-window length; the April 2020 gap;
the native 10-minute grid; and the control that retrains the raw LSTM on
exactly the windows available to the decomposition arms.}

% =====================================================================
\section{Discussion}\label{sec:discussion}
% =====================================================================
\carry{draft v2 §7, rewritten around the cost axis rather than around the raw
gap. Required changes: (i) lead with the operational consequence --- honest
evaluation is affordable; (ii) state what this adds to
Jeon et al.~\cite{jeon2026}, Zhang et al.~\cite{zhang2025} and
Zhou et al.~\cite{zhou2024} rather than restating the measurement as new;
(iii) consolidate the per-paragraph caveats into Section~\ref{sec:limits};
(iv) keep the paragraph explaining that two-sided filtering transfers more
future information than full-series VMD in these data, and connect it to the
certificate profile, which now offers a mechanism for that difference.}

\subsection{Limitations}\label{sec:limits}
\carry{draft v2 §7.1 nearly intact --- it is thorough. Add: the certificate is
a necessary condition over a probed set; $\lambda$ depends on the oracle
specification and the VMD oracle is weak; the literature search is
summary-level, and the search strings and date should be reported.}

% =====================================================================
\section{Conclusion}\label{sec:conclusion}
% =====================================================================
\carry{draft v2 §8, rewritten around the three instruments and the cost
result rather than around the raw gap.}

% =====================================================================
\section*{CRediT authorship contribution statement}
\carry{draft v2. Resolve the author-order inconsistency between the
manuscript, \texttt{CITATION.cff} and the Zenodo record before submission.}

\section*{Declaration of competing interest}
The authors declare that they have no known competing financial interests or
personal relationships that could have appeared to influence the work reported
in this paper.

\section*{Declaration of generative AI and AI-assisted technologies in the
writing process}
\carry{draft v2. Elsevier requires this statement and prescribes its wording;
check the current form in the guide for authors.}

\section*{Data availability}
The SDWPF dataset is publicly available~\cite{zhou2024sdwpf}. The code that
produces every number in this paper is available at
\url{https://github.com/GokhanGULMEZ/ruzgar-sizinti-protokolu}.
\textbf{Before submission:} the repository is currently private and this URL
does not resolve for readers; it must be made public, the stored result tables
must be added (the repository currently ships code only), and a persistent
Zenodo DOI must be minted and cited here.

% =====================================================================
\appendix
\section{Diebold--Mariano comparisons}
\carry{draft v2 Tables A1--A3.}
\section{All error metrics and robustness tables}
\carry{draft v2 Tables B1--B6, plus Tables 3, 7, 9--12, 14--19 moved here from
the main text.}
\section{CatBoost feature importance}
\carry{draft v2 Table C1 and Figures C1--C2.}

% =====================================================================
\begin{thebibliography}{99}
% Peer-reviewed sources only, except where marked. Entries carrying
% [VERIFY] need a field confirmed before submission.

\bibitem{jeon2026}
B.-K. Jeon, J. Heo, H. Joo, D. Kim, J. Kim,
Data leakage from full-series signal decomposition inflates the reported
accuracy of short-term energy forecasting hybrids: a leakage-controlled
re-evaluation,
Int. J. Energy Res. (2026).
\url{https://doi.org/10.1155/er/4433079}

\bibitem{zhang2025}
D. Zhang, K. Jiang, C. Wu, Y. Peng, Z. Han, Y. Chen,
Revisiting the efficacy of signal decomposition in AI-based time series
prediction (2025).
\url{https://doi.org/10.1117/12.3045386}

\bibitem{zhou2024}
P. Zhou, L. Shen, Y. Han, L. Mi, G. Xu,
A short-term wind speed prediction method utilizing rolling decomposition and
time-series extension to avoid information leakage,
Energy Sources Part A 46 (2024).
\url{https://doi.org/10.1080/15567036.2024.2318485}

\bibitem{valdivieso2026}
V. Valdivieso, E. López-Lao, A. García-Chica, A. Cama-Pinto,
F. M. Arrabal-Campos,
Leakage, not decomposition: auditing decompose-then-split hybrids for
short-term load forecasting, SSRN preprint (2026).
\url{https://doi.org/10.2139/ssrn.7422299}
\textsf{[PREPRINT --- replace if a peer-reviewed version appears]}

\bibitem{huo2024}
J.-Y. Huo, et al.,
Ultra-short-term wind power prediction model based on fixed scale dual mode
decomposition and deep learning networks,
Eng. Appl. Artif. Intell. 133 (2024). \textsf{[VERIFY article no./DOI]}

\bibitem{pang2026}
J. Pang, S. Dong,
Stop unrealistic data preprocessing in wind speed forecasting: approaches and
discussions on preventing future data leakage,
Appl. Math. Model. (2026). \textsf{[VERIFY volume/DOI]}

\bibitem{wangj2024}
J. Wang, W. Yuan, S. Zhang, S. Cheng, L. Han,
Implementing ultra-short-term wind power forecasting without information
leakage through cascade decomposition and attention mechanism,
Energy 312 (2024).

\bibitem{xu2023}
K. Xu, W. Wang,
Limited information limits accuracy: whether EEMD improves crude oil spot price
prediction?,
Int. Rev. Financ. Anal. 87 (2023).
\url{https://doi.org/10.1016/j.irfa.2023.102625}

\bibitem{xuA2025}
A. Xu, J. Chen, J. Li, Z. Chen, S. Xu, Y. Nie,
Multivariate rolling decomposition hybrid learning paradigm for power load
forecasting,
Renew. Sustain. Energy Rev. 212 (2025).
\url{https://doi.org/10.1016/j.rser.2025.115375}

\bibitem{yu2021}
L. Yu, Y. Ma, M. Ma,
An effective rolling decomposition-ensemble model for gasoline consumption
forecasting,
Energy 222 (2021).
\url{https://doi.org/10.1016/j.energy.2021.119869}

\bibitem{yangD2024}
D. Yang, M. Li, J. Guo, P. Du,
An attention-based multi-input LSTM with sliding window-based two-stage
decomposition for wind speed forecasting,
Appl. Energy 375 (2024).
\url{https://doi.org/10.1016/j.apenergy.2024.124057}

\bibitem{wang2024}
J. Wang, Y. Liu, Y. Li,
A parallel differential learning ensemble framework based on enhanced feature
extraction and anti-information leakage mechanism for ultra-short-term wind
speed forecast,
Appl. Energy 361 (2024) 122909.
\url{https://doi.org/10.1016/j.apenergy.2024.122909}

\bibitem{tan2025}
J. Tan, H. Zhu, J. Zhang, H. Liu,
Multi-stage wind speed prediction with CEEMDAN-SE-IDBO-LSTM based on rolling
decomposition,
Energy 338 (2025) 138741.
\url{https://doi.org/10.1016/j.energy.2025.138741}

\bibitem{chen2022}
Y. Chen, S. S. Yu, S. N. Islam, C. P. Lim, S. M. Muyeen,
Decomposition-based wind power forecasting models and their boundary issue: an
in-depth review and comprehensive discussion on potential solutions,
Energy Rep. 8 (2022) 8805--8820.
\url{https://doi.org/10.1016/j.egyr.2022.07.005}

\bibitem{yang2024}
X. Yang, J. Li, X. Jiang,
Research on information leakage in time series prediction based on empirical
mode decomposition,
Sci. Rep. 14 (2024) 28362.
\url{https://doi.org/10.1038/s41598-024-80018-9}

\bibitem{feng2026}
W. Feng, R. Tao, J. Cartlidge, J. Zheng,
VMDNet: temporal leakage-free variational mode decomposition for electricity
demand forecasting, EUSIPCO 2026, in press.

\bibitem{sesay2026}
M. M. Sesay,
Leakage-aware benchmarking of temporal data preparation for hourly energy
forecasting,
Energy Convers. Manag. X (2026).
\url{https://doi.org/10.1016/j.ecmx.2026.101791}

\bibitem{haq2026}
M. A. U. Haq, K. Ali, I. Javed, M. Fahrioglu,
Leakage safe machine learning for short term wind speed forecasting from SCADA,
ISGTA 2025, Springer Proc. Phys.

\bibitem{quilty2018}
J. Quilty, J. Adamowski,
Addressing the incorrect usage of wavelet-based hydrological and water
resources forecasting models,
J. Hydrol. (2018). \textsf{[VERIFY volume/pages/DOI]}

\bibitem{maslova2015}
I. Maslova, A. M. Ticlavilca, M. McKee,
Adjusting wavelet-based multiresolution analysis boundary conditions for
long-term streamflow forecasting,
Hydrol. Process. (2015).
\url{https://doi.org/10.1002/hyp.10564}

\bibitem{kaufman2012}
S. Kaufman, S. Rosset, C. Perlich, O. Stitelman,
Leakage in data mining: formulation, detection, and avoidance,
ACM Trans. Knowl. Discov. Data 6 (4) (2012).
\url{https://doi.org/10.1145/2382577.2382579}

\bibitem{kapoor2023}
S. Kapoor, A. Narayanan,
Leakage and the reproducibility crisis in machine-learning-based science,
Patterns 4 (9) (2023) 100804.
\url{https://doi.org/10.1016/j.patter.2023.100804}

\bibitem{rosenblatt2024}
M. Rosenblatt, et al.,
Data leakage inflates prediction performance in connectome-based machine
learning models,
Nat. Commun. 15 (2024).
\url{https://doi.org/10.1038/s41467-024-46150-w}

\bibitem{hewamalage2023}
H. Hewamalage, K. Ackermann, C. Bergmeir,
Forecast evaluation for data scientists: common pitfalls and best practices,
Data Min. Knowl. Discov. 37 (2) (2023) 788--832.
\url{https://doi.org/10.1007/s10618-022-00894-5}

\bibitem{nielsen1998}
T. S. Nielsen, A. Joensen, H. Madsen, L. Landberg, G. Giebel,
A new reference for wind power forecasting,
Wind Energy 1 (1998) 29--34.

\bibitem{madsen2005}
H. Madsen, P. Pinson, G. Kariniotakis, H. Aa. Nielsen, T. S. Nielsen,
Standardizing the performance evaluation of short-term wind power prediction
models,
Wind Eng. 29 (6) (2005) 475--489.
\url{https://doi.org/10.1260/030952405776234599}

\bibitem{messner2020}
J. W. Messner, P. Pinson, J. Browell, M. B. Bjerregård, I. Schicker,
Evaluation of wind power forecasts --- an up-to-date view,
Wind Energy 23 (6) (2020) 1461--1481.
\url{https://doi.org/10.1002/we.2497}

\bibitem{sopena2020}
J. M. González Sopeña, V. Pakrashi, B. Ghosh,
An overview of performance evaluation metrics for short-term statistical wind
power forecasting,
Renew. Sustain. Energy Rev. 138 (2020) 110515.
\url{https://doi.org/10.1016/j.rser.2020.110515}

\bibitem{sopena2022}
J. M. González Sopeña, V. Pakrashi, B. Ghosh,
A benchmarking framework for performance evaluation of statistical wind power
forecasting models,
Sustain. Energy Technol. Assess. (2022).
\url{https://doi.org/10.1016/j.seta.2023.103246}

\bibitem{lee2022}
J. C. Y. Lee, C. Draxl, L. K. Berg,
Evaluating wind speed and power forecasts for wind energy applications using an
open-source and systematic validation framework,
Renew. Energy 200 (2022) 457--475.

\bibitem{hong2020}
T. Hong, P. Pinson, Y. Wang, R. Weron, D. Yang, H. Zareipour,
Energy forecasting: a review and outlook,
IEEE Open Access J. Power Energy 7 (2020) 376--388. \textsf{[VERIFY]}

\bibitem{deignan2025}
C. Deignan, J. M. González Sopeña, B. Ghosh, V. Pakrashi,
Portability of short term wind power forecasting: investigating model
calibration using wind power data from Ireland and UK,
Sci. Rep. (2025).
\url{https://doi.org/10.1038/s41598-025-08751-3}

\bibitem{qian2019}
Z. Qian, Y. Pei, H. Zareipour, N. Chen,
A review and discussion of decomposition-based hybrid models for wind energy
forecasting applications,
Appl. Energy 235 (2019) 939--953.
\url{https://doi.org/10.1016/j.apenergy.2018.10.080}

\bibitem{dragomiretskiy2014}
K. Dragomiretskiy, D. Zosso,
Variational mode decomposition,
IEEE Trans. Signal Process. 62 (3) (2014) 531--544.
\url{https://doi.org/10.1109/TSP.2013.2288675}

\bibitem{guo2021}
Y. Guo, et al.,
Generalized variational mode decomposition: a multiscale and fixed-frequency
decomposition algorithm,
IEEE Trans. Instrum. Meas. (2021). \textsf{[VERIFY volume/article no./DOI]}

\bibitem{jia2026}
H. Jia, J. Cao, Z. Liang, et al.,
Short-time variational mode decomposition,
Signal Process. 238 (2026) 110203.
\url{https://doi.org/10.1016/j.sigpro.2025.110203}

\bibitem{dan2023}
D. Dan, et al.,
Block-wise recursive sliding variational mode decomposition method and its
application on online separating of bridge vehicle-induced strain monitoring
signals,
Mech. Syst. Signal Process. (2023). \textsf{[VERIFY volume/DOI; confirm
causality from full text before use as a baseline]}

\bibitem{wheatcroft2019}
E. Wheatcroft,
Interpreting the skill score form of forecast performance metrics,
Int. J. Forecast. (2019). \textsf{[VERIFY volume/pages/DOI]}

\bibitem{dm1995}
F. X. Diebold, R. S. Mariano,
Comparing predictive accuracy,
J. Bus. Econ. Stat. 13 (3) (1995) 253--263.

\bibitem{hln1997}
D. Harvey, S. Leybourne, P. Newbold,
Testing the equality of prediction mean squared errors,
Int. J. Forecast. 13 (2) (1997) 281--291.

\bibitem{holm1979}
S. Holm,
A simple sequentially rejective multiple test procedure,
Scand. J. Stat. 6 (2) (1979) 65--70.

\bibitem{kunsch1989}
H. R. Künsch,
The jackknife and the bootstrap for general stationary observations,
Ann. Stat. 17 (3) (1989) 1217--1241.
\url{https://doi.org/10.1214/aos/1176347265}

\bibitem{zhou2024sdwpf}
J. Zhou, X. Lu, Y. Xiao, et al.,
SDWPF: a dataset for spatial dynamic wind power forecasting over a large
turbine array,
Sci. Data 11 (2024) 649.
\url{https://doi.org/10.1038/s41597-024-03427-5}

\bibitem{kelmarsh2022}
C. Plumley,
Kelmarsh wind farm data, Zenodo (2022).
\url{https://doi.org/10.5281/zenodo.5841834}

\end{thebibliography}

\end{document}
